% ================================================================
% ==== FILE: content/cycles/cycles_k8.tex
% ================================================================

\subsubsection{Zyklen auf \texorpdfstring{$K_8$}{K_8}}
\label{sec:cycles-k8}

Das Niveau$K_8$stellt das zivilisatorische Kontinuum dar: ein großräumiges,
Mehrkomponentensystem, das Bevölkerung, Infrastrukturen, Energiebasis,
Logistik, Informationssysteme und wirtschaftliche Reproduktionsmechanismen.
Zyklen weiter$K_8$Formalisieren Sie die wiederkehrenden Makroprozesse, die a ermöglichen
Zivilisation, die sich über die historische Zeit hinweg behaupten kann.

Der Zustand von$K_8$Ist:
\[
  s(t) = (\text{Pop}(t), \text{Infra}(t), \text{Tech}(t),
           \text{Energy}(t), \text{Info}(t), \text{Econ}(t)).
\]

The fundamental flows (the corresponding derivation note) are:
\[
J^{(8)} =
  \left(
  \begin{aligned}
   &J_{\mathrm{energy}}, J_{\mathrm{material}}, J_{\mathrm{logistics}},\\
   &J_{\mathrm{info}}, J_{\mathrm{population}}
  \end{aligned}
  \right).
\]

Kritische Schwellenwerte:
\[
  \Theta_E,\;
  \Theta_L,\;
  \Theta_I,\;
  \Theta_M,\;
  \Theta_R.
\]

Strukturelle Spannung:
\[
T_8 =
  \begin{aligned}
    &w_E(E_{\mathrm{crit}} - E_{\mathrm{density}})
    + w_L(L_{\mathrm{demand}} - L_{\mathrm{capacity}})\\
    &\quad + w_M(M_{\mathrm{demand}} - M_{\mathrm{repair}})
    + w_I(I_{\mathrm{complexity}} - I_{\mathrm{coherence}})\\
    &\quad + w_R(R_{\mathrm{demand}} - R_{\mathrm{supply}})
  \end{aligned}
  .
\]

Die Existenz eines zivilisatorischen Kontinuums erfordert:
\[
  T_8 < 0,\qquad k_8>0.
\]

\subsubsection{Überblick}

Zivilisationszyklen beschreiben die makrowiederkehrenden Prozesse, durch die
Gesellschaften erhalten materielle Reproduktion, Energieflüsse, Wissenskohärenz,
Logistik, Nachhaltigkeit der Infrastruktur und demografische Stabilität.
Diese Zyklen integrieren die dynamischen Interaktionen zwischen Subsystemen und stellen sicher
the persistence of $\Omega(K_8)$.

\subsubsection{1. Energiereproduktionszyklus$C_{\mathrm{energy}}$}

Die Energiedichte ist der wichtigste Faktor für die Lebensfähigkeit einer Zivilisation.

\[
  \mathrm{Extraction}
  \to
  \mathrm{Conversion}
  \to
  \mathrm{Distribution}
  \to
  \mathrm{Use}
  \to
  \mathrm{Maintenance}
  \to
  \mathrm{Extraction}.
\]

Formal:
\[
  C_{\mathrm{energy}} =
  \left\{
    J_{\mathrm{energy}}(t),\;
    E_{\mathrm{density}}(t),\;
    \mathrm{Infra}_{\mathrm{energy}}(t)
  \right\}.
\]

Schwelle:
\[
  E_{\mathrm{density}} > \Theta_E.
\]

\subsubsection{2. Materialproduktion-Logistik-Zyklus$C_{\mathrm{mat/log}}$}

Der Makrokreislauf, der Produktion, Transport und Konsum verbindet.

\[
  \mathrm{Production}
  \to
  \mathrm{Transport}
  \to
  \mathrm{Use}
  \to
  \mathrm{Waste\;processing}
  \to
  \mathrm{Repair}
  \to
  \mathrm{Production}.
\]

Bilden:
\[
 C_{\mathrm{mat/log}} =
 \left\{
 \begin{aligned}
   &J_{\mathrm{material}}(t),\;
    J_{\mathrm{logistics}}(t),\\
   &L_{\mathrm{capacity}}(t),\;
    \mathrm{Infra}_{\mathrm{industrial}}(t)
 \end{aligned}
 \right\}.
\]

Schwelle:
\[
  L_{\mathrm{capacity}} > \Theta_L.
\]

\subsubsection{3. Wartungszyklus der Infrastruktur$C_{\mathrm{infra}}$}

Die Stabilität der Zivilisation hängt von der Aufrechterhaltung komplexer Infrastrukturen ab.

\[
  \mathrm{Inspection}
  \to
  \mathrm{Repair}
  \to
  \mathrm{Upgrade}
  \to
  \mathrm{Integration}
  \to
  \mathrm{Inspection}.
\]

Bilden:
\[
  C_{\mathrm{infra}} =
  \left\{
  \begin{aligned}
     &M_{\mathrm{repair}}(t),\;
      \mathrm{Infra}(t),\\
     &J_{\mathrm{material}}(t)
  \end{aligned}
  \right\}.
\]

Schwelle:
\[
  M_{\mathrm{repair}} > \Theta_M.
\]

\subsubsection{4. Informationskohärenzzyklus$C_{\mathrm{info}}$}

Eine Zivilisation erfordert eine konsistente Informationsverarbeitungskapazität.

\[
  \mathrm{Generation}
  \to
  \mathrm{Transmission}
  \to
  \mathrm{Integration}
  \to
  \mathrm{Coherence\;checking}
  \to
  \mathrm{Update}
  \to
  \mathrm{Generation}.
\]

Bilden:
\[
  C_{\mathrm{info}}
  =
  \left\{
  \begin{aligned}
    &J_{\mathrm{info}}(t),\;
     I_{\mathrm{coherence}}(t),\\
    &I_{\mathrm{complexity}}(t)
  \end{aligned}
  \right\}.
\]

Schwelle:
\[
  I_{\mathrm{coherence}} > \Theta_I.
\]

\subsubsection{5. Demografischer Reproduktionszyklus$C_{\mathrm{pop}}$}

Demografische Stabilität ist für die Aufrechterhaltung von Arbeitskräften, Wissenstransfer und
Militärische Kapazität und intergenerationelle Reproduktion.

\[
  \mathrm{Birth}
  \to
  \mathrm{Growth}
  \to
  \mathrm{Education}
  \to
  \mathrm{Labour}
  \to
  \mathrm{Aging}
  \to
  \mathrm{Birth}.
\]

Bilden:
\[
  C_{\mathrm{pop}}
    =
    \left\{
    \begin{aligned}
      &J_{\mathrm{population}}(t),\;
       R_{\mathrm{supply}}(t),\\
      &R_{\mathrm{demand}}(t)
    \end{aligned}
    \right\}.
\]

Schwelle:
\[
  R_{\mathrm{supply}} > \Theta_R.
\]

\subsubsection{6. Makroökonomischer Reproduktionszyklus$C_{\mathrm{econ}}$}

Der Wirtschaftskreislauf:

\[
  \mathrm{Investment}
  \to
  \mathrm{Production}
  \to
  \mathrm{Distribution}
  \to
  \mathrm{Consumption}
  \to
  \mathrm{Reinvestment}.
\]

Formal:
\[
 C_{\mathrm{econ}}=
 \left\{
 \begin{aligned}
   &\mathrm{Econ}(t),\;
    J_{\mathrm{material}}(t),\\
   &J_{\mathrm{energy}}(t),\;
    J_{\mathrm{info}}(t)
 \end{aligned}
  \right\}.
\]

Stabilität erfordert:
\[
  T_{\mathrm{econ}} < \min(\Theta_L, \Theta_E).
\]

\subsubsection{7. Integrierter Zivilisationszyklus$C_{\mathrm{civ}}$}

Das zivilisatorische Kontinuum bleibt nur bestehen, wenn alle wichtigen Reproduktionszyklen stattfinden
bleiben innerhalb stabiler Schwellenwerte.

\[
 C_{\mathrm{civ}}
   = \begin{aligned}[t]
     &C_{\mathrm{energy}}
      \cup C_{\mathrm{mat/log}}
      \cup C_{\mathrm{infra}}\\
     &\cup C_{\mathrm{info}}
      \cup C_{\mathrm{pop}}
      \cup C_{\mathrm{econ}}
   \end{aligned}
   .
\]

Stabilität erfordert:
\[
T_8
   <
   \min\!\bigl(\Theta_E,\Theta_L,\Theta_I,\Theta_M,\Theta_R\bigr).
\]

If any subthreshold fails, $\Omega(K_8)$ undergoes collapse.

\subsubsection{Zyklusmetriken für \texorpdfstring{$K_8$}{K_8}}

\paragraph{Cycle length.}
\[
  L(C) = \oint d\Omega(K_8).
\]

\paragraph{Efficiency.}
\[
  C_{\mathrm{eff}}
  = \frac{\oint J^{(8)}\cdot dA}{L(C)}.
\]

\paragraph{Stability.}
\[
SC)
    = \min_{t\in C}
      d(\Omega(K_8),\partial\Omega(K_8)).
\]

\paragraph{Weight.}
\[
WC)
 =
 F(C_{\mathrm{eff}},
SC),
   E_{\mathrm{density}},
   L_{\mathrm{capacity}},
   I_{\mathrm{coherence}}).
\]

\subsubsection{Zusammenbruch von \texorpdfstring{$K_8$}{K_8} Zyklen}

Ein Kollaps tritt auf, wenn:

\[
  E_{\mathrm{density}} < \Theta_E,\qquad
  L_{\mathrm{capacity}} < \Theta_L,\qquad
  I_{\mathrm{coherence}} < \Theta_I,\qquad
  M_{\mathrm{repair}} < \Theta_M,\qquad
  R_{\mathrm{supply}} < \Theta_R.
\]

Unter Zusammenbruch:

\[
  C_j \to \varnothing,\qquad
  \Omega(K_8)=\varnothing,\qquad
  k_8\to 0.
\]

Dies entspricht einem völligen Zusammenbruch der Zivilisation.

\subsubsection{Kontinuität von \texorpdfstring{$K_7$}{K_7} zu \texorpdfstring{$K_8$}{K_8}}

The transition $K_7\to K_8$ occurs when:

\[
  T_{\mathrm{soc}} > \Theta_{\mathrm{dim}},
\]

und neue Achsen erscheinen:

\[
  A_{\mathrm{energy}},\;
  A_{\mathrm{logistics}},\;
  A_{\mathrm{infra}},\;
  A_{\mathrm{population}},\;
  A_{\mathrm{econ}}.
\]

Dies führt zu Makroreproduktionszyklen ohne Analogie auf sozialer Ebene
$K_7$.
